\scanpage{126}
dont le corps résiduel (contenu dans celui de $V$) sera \deleted{alg} fini sur $K$ si celui de $V$ l'est (donc $V_1\ne L_1$) --- donc $V_1$ \deleted{$\longrightarrow$} \inserted{correspond à} une ``place'' du corps de fonctions d'une variable $L_1$ sur $K$. Or, par des fonctorialités essentiellement évidentes, on voit que l'hom $\mathcal E_L^{\overline L}\longrightarrow\mathcal E_{L_1}^{\overline{L_1}}$ envoie $N_V$ dans $N_{V_1}$, $I_V$ dans $I_{V_1}$ --- donc l'image de $N_V$ normalise $I_{V_1}$, donc (si on a une \inserted{germe de} section de $N_V$ sur $\mathcal E_K$) le germe de section obtenu de $\mathcal E_{L_1}$ sur $\mathcal E_K$ normalise $I_{V_1}$ --- et sa restriction à $\mathcal E_K^0$ centralise $I_{V_1}$ ---
\margin{\unsure{Il|?} faut vérifier que le corps des fractions de $V_1$ est bien $L_1$ --- c'est un fait général\ldots}

\uline{Conjecture} Soient $K,L$ des ext. de type fini de $\mathbf Q$, $K\subset L$. Alors

a) Toute section de $\mathcal E_L$ sur $\mathcal E_K$ (ou germe de tel, ça revient au $=$\ldots) normalise un $I_V$ associé à un anneau de valuation \inserted{$V$} de $L$ contenant $K$, à corps résiduel algébrique \inserted{sur $K$} \deleted{\unsure{,}} et \unsure{ce} $V$ est uniquement
\margin{125}
